\section{Conclusion}
This work designed a functionally graded alloy for the fusion first wall as a
doubly monotonic path through an eleven-element refractory composition space.
The design proceeds in four stages. Activation screening referenced to pure
tungsten and to EUROFER97 exploits the linearity of the FLARE responses, so
that the feasible space is obtained by linear programming without evaluating
the more than thirty million rejected candidates. CALPHAD-based property
filtering then enforces single-phase BCC stability and thermal conductivity,
a traversal graph is constructed on the surviving compositions, and the
shortest path is found under monotonicity imposed as a hard constraint, by
deleting every transition that raises the solidus or lowers the Pugh ratio
before an exact search. Of the
\(61{,}298\) unique activation-feasible compositions, \(700\) survive
the property screens, \(686\) of them forming a single connected design
region in which high-solidus plasma-facing candidates and ductile structural
candidates occupy opposite ends of a continuous trade-off. Between deployable
terminals, V\(_{5}\)Re\(_{5}\)W\(_{90}\) and Ti\(_{5}\)V\(_{90}\)Re\(_{5}\),
a strict doubly-monotonic path of nineteen alloys exists, a
tungsten-vanadium substitutional gradient at constant \(5\)~at.\% Re in
which the solidus only decreases, from \(3449\) to \(2195\)~K, and the Pugh
ratio only increases, from \(2.00\) to \(3.34\), at every one of its eighteen
single-exchange steps. The path length equals the minimum possible for these
terminals, so the monotonicity requirement is achieved at zero cost, and the
eighteen compositions satisfying both design criteria prove to be the sole
monotone gateway between the two categories, the path crossing three of
them. An edge-localized-mode assessment independently validates the design,
with the plasma-facing terminal carrying a melting margin of \(+1650\)~K and
the margin remaining positive along the gradient precisely through its
plasma-facing extent, the structural layers behind it sitting shielded from
surface loading.

The formulation generalizes in both of its parts. The constraint-first
screening applies to any model whose responses are linear in composition, and
the hard-constraint path planning applies to any graded-alloy design
requiring monotonicity in one or more properties, with no tuning parameters
to select. Future work will pursue validation, realization, and
generalization through experimental characterization of the
descriptor-level properties along the path, synthesis of the terminal and
gateway compositions, a graded build to test the predicted through-wall
profile, and extension of the path formulation to multi-objective design, in
which path length is balanced against additional objectives such as
thermal-expansion mismatch or activation margin.